\documentclass{article}
\usepackage[utf8]{inputenc}
\usepackage{caption}
\usepackage[margin=1in]{geometry}
\usepackage{graphicx}
\title{System Requirements Specification}
\author{Team EVAL}
\date{October 26, 2018}

\begin{document}

\maketitle

\newpage

\begin{center}
{\includegraphics[scale=.3]{team_logo.png}}\\ \bigskip
{\LARGE College Course Evaluation System }\\ \medskip
{\large System Requirements Specification }\\
\end{center}

\tableofcontents

\newpage

\section{Introduction}
\subsection{Purpose of This Document}

This system requirements specification details what our course evaluation system does and what tests we will make to ensure the system is complete. It includes why we are creating the system, the scope of the product, diagrams that illustrate the system, what we will deliver to the customer, and currently pending issues. This document is intended for the product client, Harlan Onsrud, and potential users of the system.

\subsection{References}
\subsection{Purpose of the Product}

The University of Maine gives out course evaluation surveys to students at the end of each course. The survey is filled on a bubble sheet and is then scanned. Harlan Onsrud finds it inconvenient for the school administrators to manually scan and compile the survey results. He desires an online, automated evaluation system to improve productivity.\par

With this new product, course teachers and administrators can customize their lists of questions for course evaluations and send the lists out to students. Students will receive an e-mail telling them to complete the survey, and they will be periodically notified if the survey is incomplete. After survey submissions, the teachers and administrators can view the average of the results for each question over a certain course, instructor, department, and university.

\subsection{Product Scope}

\section{Functional Requirements}

The functional requirements specify what actions the program will perform. Each requirement is represented as a use case.

\begin{center}
\vspace{4in}
\captionof{table}{}
\begin{tabular}{|p{3.5cm}|p{7.5cm}|} 
\hline
\textbf{Number} & 1  \\
\hline
\textbf{Name} & Set up course  \\ 
\hline
\textbf{Summary} & An administrator or teacher adds a course to be surveyed along with its student roll \\ 
\hline
\textbf{Priority} & 5\\ 
\hline
\textbf{Preconditions }& The course's evaluation survey is unwritten \\ 
\hline
\textbf{Postconditions} & The course's evaluation survey is written, A student roll is linked to the course \\ 
\hline
\textbf{Primary Actors }& Administrator, Teacher \\ 
\hline
\textbf{Secondary Actors} & Database \\ 
\hline
\textbf{Trigger }& A new course needs to be evaluated \\ 
\hline
\textbf{Main Scenario }& 
\begin{tabular}{l|p{5.8cm}} 
\textbf{Step }& \textbf{Action}\\
\hline
1 & Select option to add a course \\
\hline
2 & Enter the course name and \newline corresponding student roster\\
\hline
3 & Database updated with course name and roster\\
\end{tabular}\\ 
\hline
\textbf{Extensions }&
\begin{tabular}{l|p{5.8cm}} 
\textbf{Step }& \textbf{Branching Action}\\
\hline
2a & No students are entered for a course : Notify user of the problem  \\
\end{tabular}\\
\hline
\textbf{Open Issues} & What information about the course shall be stored \\ 
\hline
\end{tabular}

\bigskip
\captionof{table}{}
\begin{tabular}{|p{3.5cm}|p{7.5cm}|} 
\hline
\textbf{Number} & 2  \\
\hline
\textbf{Name} & Edit survey questions  \\ 
\hline
\textbf{Summary} & An administrator or teacher edits the survey questions for a course \\ 
\hline
\textbf{Priority} & 5\\ 
\hline
\textbf{Preconditions }& A course for the survey questions exists in the database \\ 
\hline
\textbf{Postconditions} & A course includes its survey questions \\ 
\hline
\textbf{Primary Actors }& Administrator, Teacher \\ 
\hline
\textbf{Secondary Actors} & Database \\ 
\hline
\textbf{Trigger }& A college course has zero or wrong evalution survey questions \\ 
\hline
\textbf{Main Scenario }& 
\begin{tabular}{l|p{5.8cm}} 
\textbf{Step }& \textbf{Action}\\
\hline
1 & Choose the course whose questions are to be edited \\
\hline
2 & Select questions preset if applicable\\
\hline
3 & Enter or edit the questions\\
\hline
4 & Questions stored in database and as text file\\
\end{tabular}\\ 
\hline
\textbf{Extensions }&
\begin{tabular}{l|p{5.8cm}} 
\textbf{Step }& \textbf{Branching Action}\\
\hline
3a & Invalid questions submitted : Notify user of invalid questions  \\
\end{tabular}\\
\hline
\textbf{Open Issues} & In what format shall the questions be stored in the text file \\ 
\hline
\end{tabular}

\bigskip
\vspace{2.6in}
\captionof{table}{}
\begin{tabular}{|p{3.5cm}|p{7.5cm}|} 
\hline
\textbf{Number} & 3  \\
\hline
\textbf{Name} & Create questions list preset  \\ 
\hline
\textbf{Summary} & An administrator or teacher makes a preset list of questions for a survey\\
\hline
\textbf{Priority} & 4\\ 
\hline
\textbf{Preconditions }& A course for the survey questions exists in the database \\ 
\hline
\textbf{Postconditions} & A preset list is available while creating a survey \\ 
\hline
\textbf{Primary Actors }& Administrator, Teacher \\ 
\hline
\textbf{Secondary Actors} & Database \\ 
\hline
\textbf{Trigger }& A user wants the same questions to be in multiple surveys \\ 
\hline
\textbf{Main Scenario }& 
\begin{tabular}{l|p{5.8cm}} 
\textbf{Step }& \textbf{Action}\\
\hline
1 & Select option to add a preset questions list \\
\hline
2 & Add the questions to the preset\\
\hline
3 & Database updated with preset\\
\end{tabular}\\ 
\hline
\textbf{Extensions }&
\begin{tabular}{l|p{5.8cm}} 
\textbf{Step }& \textbf{Branching Action}\\
\hline
2a & Questions in preset are invalid : Notify user of invalid questions  \\
\end{tabular}\\
\hline
\textbf{Open Issues} & Which presets are visible to each teacher \\ 
\hline
\end{tabular}

\bigskip
\captionof{table}{}
\begin{tabular}{|p{3.5cm}|p{7.5cm}|} 
\hline
\textbf{Number} & 4  \\
\hline
\textbf{Name} & Notify students  \\ 
\hline
\textbf{Summary} & A teacher e-mails students about the survey, giving automatic reminders to complete it \\ 
\hline
\textbf{Priority} & 5\\ 
\hline
\textbf{Preconditions }& A course evaluation survey is completed \\ 
\hline
\textbf{Postconditions} & All students are finished with their survey \\ 
\hline
\textbf{Primary Actors }& Teacher \\ 
\hline
\textbf{Secondary Actors} & Student, LimeSurvey \\ 
\hline
\textbf{Trigger }& The teacher releases the survey to students \\ 
\hline
\textbf{Main Scenario }& 
\begin{tabular}{l|p{5.8cm}} 
\textbf{Step }& \textbf{Action}\\
\hline
1 & Send initial reminder to students \newline about survey \\
\hline
2 & Students complete survey\\
\hline
3 & Copy LimeSurvey results to database\\
\end{tabular}\\ 
\hline
\textbf{Extensions }&
\begin{tabular}{l|p{5.8cm}} 
\textbf{Step }& \textbf{Branching Action}\\
\hline
2a & A student does not complete survey  : Send another reminder to student  \\
\end{tabular}\\
\hline
\textbf{Open Issues} & How often to remind students, handling signed comments \\ 
\hline
\end{tabular}

\bigskip
\vspace{2.6in}
\captionof{table}{}
\begin{tabular}{|p{3.5cm}|p{7.5cm}|} 
\hline
\textbf{Number} & 5 \\
\hline
\textbf{Name} & View survey results  \\ 
\hline
\textbf{Summary} & An administrator or teacher views averages of the survey results \\ 
\hline
\textbf{Priority} & 4\\ 
\hline
\textbf{Preconditions }& Survey responses are entered in database \\ 
\hline
\textbf{Postconditions} & All appropriate averages are computed and displayed \\ 
\hline
\textbf{Primary Actors }& Administrator, Teacher \\ 
\hline
\textbf{Secondary Actors} & Database \\ 
\hline
\textbf{Trigger }& A user seeks information about the responses to a group of surveys \\ 
\hline
\textbf{Main Scenario }& 
\begin{tabular}{l|p{5.8cm}} 
\textbf{Step }& \textbf{Action}\\
\hline
1 & Select option to view results \\
\hline
2 & Compute average score for each question \\
\hline
3 & Group averages by instructor, department, and university \\
\hline
4 & Display averages to user\\
\end{tabular}\\ 
\hline
\textbf{Open Issues} & Which survey questions are averaged for display \\ 
\hline
\end{tabular}
\end{center}

\subsection{Test Number(s)}

These are the tests that will verify the functional requirements:

\begin{enumerate}
  \item We will create several courses to ensure that each is loaded up with the default list of questions, and will test adding students' emails to the list of enrolled students.
  \item We will run the program several times, each time selecting different questions to be included in the survey. We will ensure some questions are mandatory and must be included no matter what.
  \item We will create a list of questions for a given course and save it. Then we will exit the program and re-login to ensure that the questions were saved.
  \item We will run several surveys with different questions and test that the averages for each question show up properly and are viewable by the administrator or teacher.
  \item After a survey is created we will enter emails and check whether the emails actually send out or not. We will also set the reminder time very low for testing purposes to see if reminder emails will actually send.

  \item (Use Case 1) The tester selects the option to add a course, then submits the course's information and student roll. The test is successful only if the course information and student roll are reflected in the database.
  \item (Use Case 2) The tester selects the option to create a survey. Next, he or she chooses a questions preset for the survey. The tester then adds one question and edits an existing question from the preset. After submitting, the database should include the new survey, containing the preset list, added question, and edited question. There must also be a new text file with the same information.
  \item (Use Case 3) The tester selects the option to create a questions preset. He or she then enters five questions into the preset and submits the preset. The test is successful only if the preset with those five questions is stored in the database.
  \item (Use Case 4) The tester messages a group of ten people about a mock survey. In a successful test, all ten people receive an e-mail containing a link to the survey. Most respondents promptly complete the survey, but one does not. The database must contain the survey responses of the nine people who finished, and the program must send another e-mail to the unfinished person about the survey.
  \item (Use Case 5) The tester selects the option to view the survey results. There are survey results for a mock instructor, mock department, and mock university. If the test is successful, the correct average scores for every question in the course evaluation survey are shown in the application. Also, the software must show the average for the mock instructor, department, and university.
>>>>>>> origin/docs
\end{enumerate}

\section{Non-Functional Requirements}

The non-functional requirements state the qualities of the program that are unrelated to its function.

\begin{center}
\captionof{table}{}
\begin{tabular}{|p{3.5cm}|p{7.5cm}|} 
\hline
\textbf{Number} & 1  \\
\hline
\textbf{Priority} & 3\\ 
\hline
\textbf{Description} & The software should be supported by the latest versions of Windows, Mac, Linux, iOS, and Android. \\ 
\hline
\textbf{Test Number(s) }& 1 \\ 
\hline
\end{tabular}

\bigskip
\captionof{table}{}
\begin{tabular}{|p{3.5cm}|p{7.5cm}|} 
\hline
\textbf{Number} & 2  \\
\hline
\textbf{Priority} & 4\\ 
\hline
\textbf{Description} & The software should be accessible by the latest versions of Safari, Chrome, Firefox, and Edge. \\ 
\hline
\textbf{Test Number(s) }& 2 \\ 
\hline
\end{tabular}


\bigskip
\captionof{table}{}
\begin{tabular}{|p{3.5cm}|p{7.5cm}|} 
\hline
\textbf{Number} & 3  \\
\hline
\textbf{Priority} & 5\\ 
\hline
\textbf{Description} & All questions entered by the teacher or administrator shall appear on the output survey. \\ 
\hline
\textbf{Test Number(s) }& 3 \\ 
\hline
\end{tabular}


\bigskip
\captionof{table}{}
\begin{tabular}{|p{3.5cm}|p{7.5cm}|} 
\hline
\textbf{Number} & 4  \\
\hline
\textbf{Priority} & 5\\ 
\hline
\textbf{Description} & All data stored in the program's database shall be valid. \\ 
\hline
\textbf{Test Number(s) }& 4 \\ 
\hline
\end{tabular}


\bigskip
\captionof{table}{}
\begin{tabular}{|p{3.5cm}|p{7.5cm}|} 
\hline
\textbf{Number} & 5  \\
\hline
\textbf{Priority} & 5\\ 
\hline
\textbf{Description} & All collected survey data shall not be alterable.\\ 
\hline
\textbf{Test Number(s) }& 5 \\ 
\hline
\end{tabular}


\bigskip
\captionof{table}{}
\begin{tabular}{|p{3.5cm}|p{7.5cm}|} 
\hline
\textbf{Number} & 6  \\
\hline
\textbf{Priority} & 4\\ 
\hline
\textbf{Description} & Teachers shall not be able to access data of courses other than their own. \\ 
\hline
\textbf{Test Number(s) }& 6 \\ 
\hline
\end{tabular}


\bigskip
\captionof{table}{}
\begin{tabular}{|p{3.5cm}|p{7.5cm}|} 
\hline
\textbf{Number} & 7  \\
\hline
\textbf{Priority} & 3\\ 
\hline
\textbf{Description} & The mean time between failures should be at least 60 minutes. \\ 
\hline
\textbf{Test Number(s) }& 7 \\ 
\hline
\end{tabular}


\bigskip
\captionof{table}{}
\begin{tabular}{|p{3.5cm}|p{7.5cm}|} 
\hline
\textbf{Number} & 8  \\
\hline
\textbf{Priority} & 5\\ 
\hline
\textbf{Description} & Students shall have no access to any data stored by the program. \\ 
\hline
\textbf{Test Number(s) }& 8 \\ 
\hline
\end{tabular}


\bigskip
\captionof{table}{}
\begin{tabular}{|p{3.5cm}|p{7.5cm}|} 
\hline
\textbf{Number} & 9  \\
\hline
\textbf{Priority} & 5 \\ 
\hline
\textbf{Description} & All survey responses shall be anonymous. \\ 
\hline
\textbf{Test Number(s) }& 9 \\ 
\hline
\end{tabular}


\bigskip
\captionof{table}{}
\begin{tabular}{|p{3.5cm}|p{7.5cm}|} 
\hline
\textbf{Number} & 10  \\
\hline
\textbf{Priority} & 2\\ 
\hline
\textbf{Description} & The software should scale to at least three universites, 1000 courses per semester, 1000 teachers per university, and 500 students per course. \\ 
\hline
\textbf{Test Number(s) }& 10 \\ 
\hline
\end{tabular}


\bigskip
\captionof{table}{}
\begin{tabular}{|p{3.5cm}|p{7.5cm}|} 
\hline
\textbf{Number} & 11  \\
\hline
\textbf{Priority} & 1 \\ 
\hline
\textbf{Description} & The software should not exceed 500 MB in size. \\ 
\hline
\textbf{Test Number(s) }& 11 \\ 
\hline
\end{tabular}


\bigskip
\captionof{table}{}
\begin{tabular}{|p{3.5cm}|p{7.5cm}|} 
\hline
\textbf{Number} & 12  \\
\hline
\textbf{Priority} & 4\\ 
\hline
\textbf{Description} & The software's source code shall be open-source. \\ 
\hline
\textbf{Test Number(s) }& 12\\ 
\hline
\end{tabular}


\bigskip
\captionof{table}{}
\begin{tabular}{|p{3.5cm}|p{7.5cm}|} 
\hline
\textbf{Number} & 13  \\
\hline
\textbf{Priority} & 4 \\ 
\hline
\textbf{Description} & The licensing requirements of any non-original code shall be met.\\ 
\hline
\textbf{Test Number(s) }& 13 \\ 
\hline
\end{tabular}

\bigskip
\vspace{.5in}
\captionof{table}{}
\begin{tabular}{|p{3.5cm}|p{7.5cm}|} 
\hline
\textbf{Number} & 14  \\
\hline
\textbf{Priority} & 4 \\ 
\hline
\textbf{Description} & The software shall meet UMaine AFUM requirements.\\ 
\hline
\textbf{Test Number(s) }& 13 \\ 
\hline
\end{tabular}

\end{center}

\subsection{Test Number(s)}

These are the tests that will verify the non-functional requirements:

\begin{enumerate}
  \item The tester attempts to run the program on the last versions of Windows, Mac, Linux, iOS, and Android. The test is successful only if all five use cases execute properly on each operating system.
  \item The tester attempts to run the program on the last versions of Safari, Chrome, Firefox, and Edge. The test is successful only if all five use cases execute properly on each browser.
  \item The tester inputs five diverse questions into a mock survey and sends the survey to a mock student. The test is successful only if the mock student can follow the link to a survey showing those five questions.
  \item The tester inputs a course, questions preset, and survey. The test is successful only if all data for the course, preset, and survey are present in the database, and if that data correctly follows the database schema.
  \item The tester inputs the same course, preset, and survey as in test 4. The test is successful only if any function in the software does not change the survey responses stored in the database.
  \item The tester creates three users for the evaluation program. Each user inputs a different course, questions preset, and survey. The test is successful if the users cannot view nor modify each other's data.
  \item The tester runs an automatic script that simulates several hours of using the software. The test is successful only if the average time between two errors in the software is at least 60 simulated minutes.
  \item The tester inputs a course, present, and survey as a teacher, then poses as someone unauthorized to use the software. In a successful test, the tester cannot view nor modify any data present in the database.
  \item Five testers complete a survey sent by the software, and a sixth tester uses the software afterwards. The test is successful only if the sixth tester cannot find the identities of anyone who finished the survey in the database.
  \item The tester runs an automatic script that sets up the software with three colleges, 1000 courses per college, 1000 teachers per university, and 500 students per course. The test is successful if all five use cases execute properly for all courses.
  \item The tester sees the info for the folder containing the code. The folder size is at most 500 megabytes if the test passes.
  \item Ten computer users attempt to access the source code on GitHub. The test is successful if all users can do so.
  \item The tester reads the GNU General Public License and University of Maine AFUM requirements. The test is successful if the software meets all terms in both the license and AFUM requirements.
\end{enumerate}

\section{User Interface}

See ``User Interface Design Document for the College Course Evaluation System.''

\section{Deliverables}

The following lists the estimated date and format that each submission will be delivered:

\begin{center}
\captionof{table}{}
\begin{tabular}{|p{6cm}|p{3cm}|p{3cm}|} 
\hline
\textbf{Submission} & \textbf{Date of Delivery} & \textbf{Format} \\
\hline
System Requirements Specification & 10/26/2018 & Hard Copy\\ 
\hline
System Design Document & 11/16/2018 & Hard Copy\\ 
\hline
User Interface Design Document & 11/30/2018 & Hard Copy\\ 
\hline
Administrator Manual & April 2019 & Hard Copy\\ 
\hline
User Manual & May 2019 & Hard Copy\\ 
\hline
System Requirements Specification & May 2019 & Electronic\\ 
\hline
System Design Document & May 2019 & Electronic\\ 
\hline
User Interface Design Document & May 2019 & Electronic\\ 
\hline
Administrator Manual & May 2019 & Electronic\\ 
\hline
User Manual & May 2019 & Electronic\\ 
\hline
Source Code & May 2019 & Electronic\\ 
\hline
Web Link to Program & May 2019 & Electronic\\ 
\hline
\end{tabular}
\end{center}

\section{Open Issues}

The most significant issue our team has relates to the non-anonymous student responses. In a course evaluation survey, a student has the option to add a signed comment to be stored in an instructor's file. However, the LimeSurvey software does not store the identities of survey respondents. We need to find a way to collect the signed comment along with the survey, while ensuring that the signature is authentic.

There are several more minor issues that need to be resolved. We have not decided what information about each course should be stored in the database. We do not know what format to use for the questions file and how to import it into LimeSurvey. We should also know which question presets and which surveys' results are visible to each user. Finally, we should decide how often to notify students about their survey.

\appendix

\newpage
\section{Agreement Between Customer and Contractor}
This page shows that all members of Team EVAL and the client, Harlan Onsrud, have agreed on all the information in the system requirements specification. By signing this document, Team EVAL and Dr. Onsrud agree on the goals and scope of the project, each use case and requirement, the tests for each requirement, and how the deliverables will be sent.

The team will follow a process in the case that the requirements specification is changed after we sign it. ENTER THIS PROCESS HERE.

\vspace{.7in}
\noindent
\begin{tabular}{ p{5cm} p{5cm} p{5cm} } 
\textbf{\textit{Name}} & \textbf{\textit{Signature}} & \textbf{\textit{Date}} \\[.5cm]
\textbf{Jovon Craig} & $\rule{5cm}{.1mm}$ & $\rule{5cm}{.1mm}$\\[.5cm]
\textbf{Sam Elliot} & $\rule{5cm}{.1mm}$ & $\rule{5cm}{.1mm}$\\[.5cm]
\textbf{Robert Judkins} & $\rule{5cm}{.1mm}$ & $\rule{5cm}{.1mm}$\\[.5cm]
\textbf{Stanley Small} & $\rule{5cm}{.1mm}$ & $\rule{5cm}{.1mm}$\\[.5cm]
\textbf{Harlan Onsrud} & $\rule{5cm}{.1mm}$ & $\rule{5cm}{.1mm}$\\[.5cm]
Customer Comments: & \multicolumn{2}{ l }{ $\rule{10.5cm}{.1mm}$ }\\[.5cm]
\multicolumn{3}{ l }{ $\rule{16cm}{.1mm}$ }\\[.5cm]
\end{tabular}

\newpage
\section{Team Review Sign-off}

This page shows that all members of Team EVAL have reviewed the system requirements specification and agreed on its content. By signing this document, the team members agree on the goals and scope of the project, each use case and requirement, the tests for each requirement, and how the deliverables will be sent. There is nothing in the document that is a source of contention.

\vspace{.7in}
\noindent
\begin{tabular}{ p{5cm} p{5cm} p{5cm} } 
\textbf{\textit{Name}} & \textbf{\textit{Signature}} & \textbf{\textit{Date}} \\[.5cm]
\textbf{Jovon Craig} & $\rule{5cm}{.1mm}$ & $\rule{5cm}{.1mm}$\\[.5cm]
Comments: & \multicolumn{2}{ l }{ $\rule{10.5cm}{.1mm}$ }\\[.5cm]
\multicolumn{3}{ l }{ $\rule{16cm}{.1mm}$ }\\[.5cm]
\textbf{Sam Elliot} & $\rule{5cm}{.1mm}$ & $\rule{5cm}{.1mm}$\\[.5cm]
Comments: & \multicolumn{2}{ l }{ $\rule{10.5cm}{.1mm}$ }\\[.5cm]
\multicolumn{3}{ l }{ $\rule{16cm}{.1mm}$ }\\[.5cm]
\textbf{Robert Judkins} & $\rule{5cm}{.1mm}$ & $\rule{5cm}{.1mm}$\\[.5cm]
Comments: & \multicolumn{2}{ l }{ $\rule{10.5cm}{.1mm}$ }\\[.5cm]
\multicolumn{3}{ l }{ $\rule{16cm}{.1mm}$ }\\[.5cm]
\textbf{Stanley Small} & $\rule{5cm}{.1mm}$ & $\rule{5cm}{.1mm}$\\[.5cm]
Comments: & \multicolumn{2}{ l }{ $\rule{10.5cm}{.1mm}$ }\\[.5cm]
\multicolumn{3}{ l }{ $\rule{16cm}{.1mm}$ }\\[.5cm]
\end{tabular}


\newpage
\section{Document Contributions}

Stanley Small created the document, added the system requirements specification template to the document, and formatted the text such that it is compatible with LaTeX. He also ENTER WORK HERE. Stan contributed about ?? percent to the document.

Jovon Craig wrote a draft of the functional and non-functional requirements, as well as a draft of the tests. He also wrote the purpose of the document, purpose of the product, deliverables section, and the text and blank lines in the appendices. Jovon contributed about ?? percent to the document.

Sam Elliot ENTER WORK HERE. He contributed about ?? percent to the document.

Robert Judkins helped forumlate the tests for each of the functional requirements. He contributed about ?? percent to the document.


\end{document}
